\subsection{Methodology Overview}

Chapter 1 introduced the physical and geological background of thermoelastic wave propagation in geological media. It explained how rock type, mineral composition, porosity, fractures, thermal properties, elastic stiffness, and anisotropy influence the propagation of thermal and mechanical disturbances. Chapter 2 then translated these physical concepts into a mathematical formulation by defining the spatial domain, time interval, coordinate system, physical fields, material parameters, governing equations, initial and boundary conditions, and directional quantities.

The purpose of this chapter is to define the methodology used to transform the mathematical formulation into a directional prediction framework. The methodology does not replace the thermoelastic equations by an unrelated data-driven procedure. Instead, it uses the variables and parameters defined in Chapter 2 as the physical basis for the prediction task. The temperature field, displacement field, material parameters, source location, observation location, and propagation direction form the core structure of the input-output description.

The study is positioned as a research prototype for AI-assisted directional prediction of thermoelastic wave behavior in geological media. This means that the work does not claim to provide a complete field-validated simulator of real two-dimensional geological wave propagation. The full coupled thermoelastic problem in heterogeneous, fractured, anisotropic, and fluid-bearing geological media is substantially more complex than the simplified formulation used here. Therefore, the methodology focuses on a physically motivated prediction framework rather than on complete deterministic numerical simulation of a field-scale geological system.

The central methodological idea is to compare alternative predictive components within a common input-output structure. The same physical variables and material descriptions are used to represent the thermoelastic state of the medium, while different model families are treated as alternative ways of estimating the directional response. This makes the comparison more systematic: the models are not considered as independent black boxes with unrelated inputs, but as components evaluated within the same physical and mathematical context.

The methodology is therefore organized around four elements. First, the thermoelastic formulation from Chapter 2 is converted into a prediction task. Second, the input and output variables are defined in a consistent way. Third, several model families are compared as alternative predictive components. Fourth, the results are interpreted using criteria such as physical consistency, directional behavior, stability, comparison between rock types, comparison between models, interpretability, and practical availability in the prototype.
